\section{Two labelled cycles with the same ordinary spectrum}
Attach, in cyclic order, the four distinct quadratic characters of
fundamental discriminants $-3,-4,5,8$ to four vertices. The first two
characters are odd and the last two even. Start with conductances
$1,2,1,2$ around the cycle and then interchange $1$ and $2$ on every
edge. The two Laplacians $A,B$ are conjugate by a cyclic permutation, so
their ordinary eigenvalues agree with multiplicities. Their labelled
determinant polynomials satisfy the identity
\begin{equation}\label{eq:cycleexample}
 p_A(y)-p_B(y)=3(y_1-y_3)(y_2-y_4).
\end{equation}
Indeed all four diagonal entries equal $3$. For each edge, the complementary
quadratic coefficient is $9-w^2$, and for each nonedge it is $9$. All
one-variable coefficients agree because the four three-vertex principal
minors are equal for a Laplacian of a connected graph (or by direct
evaluation); both full determinants vanish. This proves
\eqref{eq:cycleexample}. The file \texttt{verify\_cycle.py} independently
checks all sixteen coefficients with integer arithmetic.

A common scalar translation of the four $y_i$ leaves the difference in
\eqref{eq:cycleexample} unchanged. As $p_B(u\mathbf1+y)$ is monic of degree
four in $u$, the first nonzero term of the log quotient has order $u^{-4}$
and coefficient $3(y_1-y_3)(y_2-y_4)$. Thus the fourth-order case of
Theorem~\ref{thm:couplinggap} actually occurs for connected graphs with
strictly positive edge weights. Theorem~\ref{thm:localcoupled} gives a
linearly long alternating sequence of excess multiplicities in each
eligible window, even in this fourth-order case. If the character labels
are rotated together with the graph, the inputs instead represent the
same labelled object up to simultaneous permutation.
